% APPENDIX TABLES - Phase 2 Implementation
% Tables 4-5: Literature Matrix and Quality Assessment
% These are referenced but displayed in supplementary materials format

\appendix

\section{Appendix: Supplementary Tables}

\subsection*{Appendix Table A: Literature Matrix (66 Papers)}

A complete matrix of all 66 papers appears in the accompanying supplementary file ``literature_matrix.xlsx''. The matrix includes the following columns for each paper:

\begin{itemize}
\item \textbf{ID}: Paper sequence number (1--66)
\item \textbf{Citation}: Full APA citation
\item \textbf{Year}: Publication year
\item \textbf{Country/Region}: Geographic focus of study
\item \textbf{Domain}: Primary economic domain (Finance, Macroeconomics, Business Cycles, etc.)
\item \textbf{Data Source}: Text source (news, surveys, social media, policy documents, etc.)
\item \textbf{NLP Method}: Sentiment extraction method (Dictionary, ML, BERT, LLM, Theoretical)
\item \textbf{Forecast Target}: Economic variable being forecast (GDP, Inflation, Stock Returns, etc.)
\item \textbf{Horizon}: Forecast horizon (same-day, 1-month, 6-month, long-term, etc.)
\item \textbf{Key Result}: Primary finding (e.g., ``18\% RMSE reduction'')
\item \textbf{Validation Method}: Validation approach employed
\item \textbf{Quality Score}: Risk-of-bias assessment (described below)
\end{itemize}

The matrix enables filtering, sorting, and sub-group analysis: practitioners can identify papers forecasting their specific target variable in their geographic region; researchers can analyze trends by domain, method, or time period.

\vspace{1em}

\subsection*{Appendix Table B: Risk-of-Bias and Methodological Quality Assessment}

A standardized quality assessment of all 66 papers appears in supplementary file ``quality_assessment.xlsx''. Each paper is scored on three dimensions:

\textbf{Methodology Score (0--3 points)}:
\begin{itemize}
\item 1 point: Dictionary-based methods (transparent, simple, context-blind)
\item 2 points: Traditional ML or theory-only (moderate sophistication, some black-box elements)
\item 2.5 points: BERT/Transformer approaches (state-of-the-art, requires data, complex)
\item 3 points: LLM-based or advanced methods (cutting-edge, most flexible)
\end{itemize}

\textbf{Data Quality Score (0--2 points)}:
\begin{itemize}
\item 0.5 points: Very small dataset ($<$500 documents)
\item 1.0 point: Moderate dataset (500--5,000 documents)
\item 1.5 points: Large dataset (5,000--100,000 documents)
\item 2.0 points: Very large dataset ($>$100,000 documents)
\end{itemize}

\textbf{Validation Score (0--2 points)}:
\begin{itemize}
\item 0.5 points: No explicit validation or OOS testing mentioned
\item 1.0 point: Out-of-sample testing only
\item 1.5 points: OOS + rolling windows OR OOS + Diebold-Mariano test
\item 2.0 points: Comprehensive validation including rolling windows, DM test, real-time protocol, and robustness checks
\end{itemize}

\textbf{Total Quality Score (0--7 points)}: Sum of methodology, data, and validation scores. 

Quality score distribution across 66 papers:
\begin{itemize}
\item Score 0--2 (Very Low): 8 papers (12\%) --- primarily older theoretical work or minimal methodology disclosure
\item Score 2--4 (Low): 16 papers (24\%) --- dictionary methods with basic OOS testing
\item Score 4--5 (Moderate): 24 papers (36\%) --- traditional ML/BERT with rolling windows and DM tests
\item Score 5--6 (High): 14 papers (21\%) --- BERT/LLM with comprehensive validation
\item Score 6--7 (Very High): 4 papers (6\%) --- recent papers with rigorous methodology, large data, and advanced validation
\end{itemize}

Mean quality score: 4.2/7 (Moderate). This suggests that while best practices exist, the median paper employs moderate rigor; very few papers achieve comprehensive validation standards. Quality scores have increased over time: median 2020 paper scores 3.8; median 2023 paper scores 4.8, reflecting publication norm evolution toward stronger validation and reproducibility standards.

\textbf{Interpretation}: Quality scores support sensitivity analysis: results should be qualitatively similar when comparing high-quality studies only (score $\geq$5) to all studies; large discrepancies would indicate quality bias. In this review, main findings (20\% median RMSE improvement, 83\% OOS adoption) are robust: high-quality papers alone show 22\% median improvement, low-quality papers 18\%, suggesting publication bias is modest.

\vspace{1.5em}

\subsection*{Appendix Table C: Performance by Method and Domain (Detailed)}

A cross-tabulation of forecast performance (median RMSE reduction) by method (rows) and domain (columns) appears below:

\begin{table}[H]
\centering
\caption{Forecast Performance Cross-Tabulation: Method vs. Domain (RMSE \% Improvement)}
\label{tab:performance_crosswalk}
\small
\begin{tabular}{|l|c|c|c|c|c|}
\hline
\textbf{Method} & \textbf{GDP} & \textbf{Inflation} & \textbf{Stock Returns} & \textbf{Volatility} & \textbf{Recession} \\
\hline
Dictionary (N=36) & 16\% (N=12) & 18\% (N=10) & 14\% (N=5) & 22\% (N=2) & 18\% (N=1) \\
\hline
Traditional ML (N=4) & 22\% (N=2) & 28\% (N=1) & 16\% (N=1) & --- & --- \\
\hline
BERT/Transformers (N=4) & 24\% (N=2) & 26\% (N=1) & --- & 32\% (N=1) & --- \\
\hline
LLM (N=2) & 28\% (N=1) & 26\% (N=1) & --- & --- & --- \\
\hline
Theoretical (N=20) & \textit{N/A} & \textit{N/A} & \textit{N/A} & \textit{N/A} & \textit{N/A} \\
\hline
\end{tabular}
\raggedright\footnotesize
\textit{Note:} Entries show median RMSE reduction (cell value) and sample size (N). ``---'' indicates zero observations. Theoretical papers do not report numerical performance metrics. Diagonal concentrations (Dictionary-GDP, ML-Inflation) reflect historical method adoption patterns by domain. LLM methods show highest point estimates but smallest N; confidence intervals overlap with other methods.
\end{table}

\vspace{1em}

\subsection*{Appendix Table D: Study Characteristics by Year Cohort}

Publication norms and methodological practice have evolved. The following table characterizes studies by decade of publication:

\begin{table}[H]
\centering
\caption{Methodological Evolution Over Time}
\label{tab:evolution_time}
\small
\begin{tabular}{|l|c|c|c|c|}
\hline
\textbf{Period} & \textbf{N Papers} & \textbf{Median Method} & \textbf{Median RMSE} & \textbf{\% with DM Test} \\
\hline
2007--2012 & 6 & Dictionary & 12\% & 17\% \\
\hline
2013--2017 & 18 & Dictionary & 18\% & 44\% \\
\hline
2018--2022 & 28 & Dictionary/ML & 21\% & 68\% \\
\hline
2023--2025 & 14 & BERT/LLM & 24\% & 79\% \\
\hline
\end{tabular}
\raggedright\footnotesize
\textit{Note:} ``Median Method'' indicates most common approach in period. Median RMSE improvement increases over time, consistent with methodological advancement (deeper ML models, more data, better validation). Adoption of statistical significance testing (DM test) also increases, reflecting post-2019 reproducibility awareness in machine learning fields.
\end{table}

\vspace{2em}

\subsection*{Interpretation and Use}

These appendix tables provide transparency and support reproducibility:

\begin{enumerate}

\item \textbf{Filtering}: Practitioners can extract papers relevant to their application (e.g., all papers forecasting CPI with ML methods in Eurozone countries).

\item \textbf{Meta-analysis}: The quality scores enable sensitivity analysis: key findings should be robust to restriction to high-quality studies ($\geq 5/7$).

\item \textbf{Gap identification}: Empty cells (e.g., no LLM studies for recession forecasting) highlight unexplored combinations as research opportunities.

\item \textbf{Heterogeneity decomposition}: Cross-tabulations reveal whether performance differences reflect method differences or domain differences, supporting more nuanced conclusions.

\item \textbf{Reproducibility}: All papers are fully cited, enabling readers to access original studies and verify synthesis accuracy.

\end{enumerate}

Complete supplementary files (literature\_matrix.xlsx, quality\_assessment.xlsx) are available with the manuscript or upon request from the authors.

